% Volume VIII: Applications and Case Structures
% Part XV. Domains of Difficult Judgment
% Chapter 89: Scientific Misconduct
% Spine: proof-spines/ch089-spine.md

\chapter{Scientific Misconduct}
\label{ch:ch089}

Scientific misconduct is a domain in which evidence about the world cannot be separated from evidence about how the result was produced. This chapter takes up authorship disputes, data irregularities, peer review, replication failure, and institutional conflict in order to distinguish anomaly, suspicion, investigation, finding, and sanction. The distinction matters because communities often collapse these stages, treating every irregularity either as scandal already proved or as a problem too embarrassing to examine.

The chapter therefore asks how whistleblowing, image screening, data access, replication, and formal inquiry should interact when the object of suspicion is not a single proposition but an evidentiary process. Its aim is to show that disciplined paranoia is internal to science itself: trust in results depends on institutions that can investigate their own production conditions without confusing preliminary concern with final condemnation.
